\begin{frame}{Outline}
  \tableofcontents
\end{frame}

% ==================================================================
\section{Motivation \& the toy model}
\begin{frame}{The problem: tracking is a combinatorial needle-in-a-haystack}
  \begin{columns}[T]
    \begin{column}{0.54\textwidth}
      \begin{itemize}
        \item The LHCb \textbf{VELO} reconstructs charged particles (\textbf{tracks}) from \textbf{hits} on silicon planes.
        \item HL-LHC pile-up $\Rightarrow$ reconstruction cost threatens to outpace classical resources.
        \item With $T$ hits per plane each plane-gap offers $T^2$ candidate links, only $T$ real:
        \[ n_{\rm seg}=4T^2 \quad\text{vs}\quad n_{\rm true}=4T \quad(\text{signal fraction } 1/T).\]
        \item At $T{=}400$: \textbf{640{,}000} candidates, \textbf{1{,}600} true (0.25\%).
        \item \textbf{Idea:} cast track finding as a \emph{linear system} $A\mathbf{x}=\mathbf{b}$, solvable classically \emph{or} on a quantum computer.
      \end{itemize}
    \end{column}
    \begin{column}{0.46\textwidth}
      \centering\fig{\textwidth}{\PF/event_display_3d.png}
      \cn{A generated toy VELO event: tracks fan from the PV through 5 planes.}
    \end{column}
  \end{columns}
\end{frame}

\begin{frame}{A controlled toy VELO --- our testbed}
  \begin{columns}[T]
    \begin{column}{0.5\textwidth}
      \begin{itemize}
        \item \texttt{LHCb\_VeLo\_Toy\_Model}: \textbf{5 planes} at $z{=}33,66,99,132,165$\,mm, $\pm40$\,mm.
        \item Deterministic and fully controlled: sweep multiplicity $T$, noise, density, kernel.
        \item Decoupled pipeline: events stored once; the matrix $A$ regenerated on demand (sparse); metrics recomputed as a \emph{view}.
        \item Purpose: isolate \emph{how each physical factor} moves the algorithm.
      \end{itemize}
    \end{column}
    \begin{column}{0.5\textwidth}
      \centering\fig{\textwidth}{\PF/pipeline_architecture.png}
    \end{column}
  \end{columns}
\end{frame}

\begin{frame}{The toy injects the two physical detector noises}
  \centering
  \figw{0.9\textwidth}{\PF/noise_types.png}
  \begin{itemize}
    \item \textbf{Multiple scattering} $\sigma_{\rm scatt}$: kicks the slopes $(t_x,t_y)$ at each plane $\Rightarrow$ a cone around the ideal track.
    \item \textbf{Resolution} $\sigma_{\rm res}$: smears each measured hit in $x,y$. \quad Both \emph{widen the triplet kink angle} $\Rightarrow$ they set $\varepsilon$.
  \end{itemize}
\end{frame}

% ==================================================================
\section{From hits to the Hamiltonian to a matrix inversion}
\begin{frame}{Step 1: from hits to segments}
  \centering
  \figw{0.93\textwidth}{\SL/walkthrough_400/event_and_segments.png}
  \begin{itemize}
    \item A \textbf{segment} is a candidate hit$\to$hit link between adjacent planes; assign each a variable $x_i$ (its \emph{activation}).
    \item \textcolor{grn}{\textbf{Green}} $=$ true (both hits from the same particle); \textcolor{acc}{\textbf{red dashed}} $=$ false (cross-track accident).
    \item The reconstruction task: switch ON the green, OFF the red --- $4T^2$ binary decisions.
  \end{itemize}
\end{frame}

\begin{frame}{Step 2: the kink-angle cost --- which segments support each other}
  \begin{columns}[T]
    \begin{column}{0.52\textwidth}
      \begin{itemize}
        \item Two segments sharing a middle hit form a \textbf{triplet} with kink angle
        $\theta_{ij}=\arccos(\hat v_i\cdot\hat v_j)$.
        \item A real particle goes nearly straight $\Rightarrow$ small $\theta$. The \textbf{cost/compatibility kernel} is explicit:
        \[ C(\theta_{ij})=\begin{cases}1 & \text{share a hit and } \theta_{ij}<\varepsilon\\[2pt] 0 & \text{otherwise}\end{cases} \]
        \item smooth variant: $C=1+\mathrm{erf}\!\big(\tfrac{\varepsilon-\theta}{\theta_d\sqrt2}\big)\in[0,2]$.
        \item $\varepsilon$ is \emph{derived from the detector noise}:
        \[ \varepsilon=\sqrt{2(s\sigma_{\rm scatt})^2+12\arctan^2\!\tfrac{s\sigma_{\rm res}}{\Delta z}+2\theta_{\min}^2}. \]
      \end{itemize}
    \end{column}
    \begin{column}{0.48\textwidth}
      \centering\fig{\textwidth}{\PF/kernel_C_theta.png}
    \end{column}
  \end{columns}
\end{frame}

\begin{frame}{Step 3: the explicit Hamiltonian}
  Reward mutually-compatible segments, penalise activation, bias every segment on:
  \[ \boxed{\;H(\mathbf{x}) \;=\; \underbrace{\frac{\gamma+\delta}{2}\sum_i x_i^2}_{\text{activation penalty}}
      \;-\; \underbrace{\frac12\sum_{i\neq j} C(\theta_{ij})\,x_i x_j}_{\text{triplet reward (attractive)}}
      \;-\; \underbrace{\delta\sum_i x_i}_{\text{bias}}\;} \]
  \begin{itemize}
    \item An Ising/Hopfield-type energy over segment activations (Denby--Peterson lineage); defaults $\gamma{=}3$, $\delta{=}1$.
    \item Physically: each active segment ``pulls up'' the segments that continue it through a shared hit at a small kink angle; lone segments are pulled down.
  \end{itemize}
\end{frame}

\begin{frame}{Step 4: minimising $H$ \emph{is} a matrix inversion}
  Stationarity, $\partial H/\partial x_i = 0$ for every segment $i$:
  \[ (\gamma+\delta)\,x_i \;-\; \sum_{j} C(\theta_{ij})\,x_j \;-\; \delta = 0
     \qquad\Longleftrightarrow\qquad \boxed{\,A\,\mathbf{x}=\delta\mathbf{1},\quad A=(\gamma+\delta)I - C\,} \]
  \begin{columns}[T]
    \begin{column}{0.5\textwidth}
      \begin{itemize}
        \item Reconstruction $=$ \textbf{solve a sparse linear system}; classically: $\mathbf{x}=A^{-1}\delta\mathbf{1}$ (LU/CG).
        \item \textbf{Hopfield fixed points} give an absolute scale:
        false/isolated $\to\delta/(\delta{+}\gamma)=0.25$;\quad
        true chain $\to\big[\tfrac{4}{11},\tfrac{5}{11}\big]$;\quad
        threshold $\tau{=}0.35$.
        \item Classical solve $=$ state-of-the-art efficiency \cn{[arXiv:2511.11458]}.
      \end{itemize}
    \end{column}
    \begin{column}{0.5\textwidth}
      \centering\figh{0.52\textheight}{\PF/activation_hopfield.png}
    \end{column}
  \end{columns}
\end{frame}

\begin{frame}{What the inversion produces on one clean 400-track event}
  \centering
  \figw{0.78\textwidth}{\PF/activation_spectrum_400.png}
\end{frame}

% ==================================================================
\section{Quantum inversion: HHL $\to$ one bit $\to$ 1BQF}
\begin{frame}{HHL: the quantum linear-system algorithm}
  \begin{columns}[T]
    \begin{column}{0.56\textwidth}
      \textbf{Harrow--Hassidim--Lloyd} solves $A\mathbf{x}=\mathbf{b}$:
      \begin{enumerate}
        \item prepare $|\mathbf{b}\rangle$ over the segment basis;
        \item \textbf{QPE}: $e^{iAt}$ writes each eigenvalue $\lambda_k$ onto an $n_{\rm time}$-qubit clock register;
        \item controlled $R_y$: amplitude $\propto 1/\lambda_k$ onto an ancilla;
        \item uncompute, post-select ancilla $=1$ $\Rightarrow$ state $\propto A^{-1}\mathbf{b}$.
      \end{enumerate}
      \medskip
      Exponential speedup in the number of hits --- \emph{contingent on} efficient QPE \& readout.\\[4pt]
      \textbf{Catch:} dense $e^{iAt}$ + high-precision QPE $\Rightarrow$ prohibitive depth on NISQ hardware.
    \end{column}
    \begin{column}{0.44\textwidth}
      \centering\fig{\textwidth}{\PF/hhl_vs_1bqf_circuit.png}
      \cn{Left: the HHL circuit. Right: where we are heading.}
    \end{column}
  \end{columns}
\end{frame}

\begin{frame}{HHL stage by stage: the statevector at every step}
  \centering
  \figw{0.99\textwidth}{\PF/hhl_states.png}
  \vspace{2pt}
  {\footnotesize On a true-track block ($P_4$): $|b\rangle \xrightarrow{\rm QPE} \sum_k\beta_k|u_k\rangle|\lambda_k\rangle
   \xrightarrow{R_y(1/\lambda)} \text{ancilla}=1 \text{ branch} \xrightarrow{\rm post\text{-}select} \;\propto A^{-1}\mathbf{b}$.}
\end{frame}

\begin{frame}{TrackHHL: restrict the phase estimation to \emph{one bit}}
  \cn{[arXiv:2511.11458 --- Chiotopoulos, Lucio Martinez, Nicotra \emph{et al.}]}
  \begin{itemize}
    \item For this Hamiltonian the spectrum is \textbf{bimodal}: a huge bulk at $\lambda=\gamma+\delta$ (lone/false segments) and a narrow band of track modes away from it.
    \item One bit of eigenvalue information --- \emph{on the bulk or off it} --- is already the answer.
    \item Restricting QPE to \textbf{one clock qubit}:
      cuts circuit depth by up to $\mathbf{10^4}$; \quad
      resolves the read-out problem (the output is binary); \quad
      enables primary-vertex post-processing.
  \end{itemize}
  \centering\figh{0.4\textheight}{\SL/single_event_400/spectrum_notch.png}\\
  \cn{The event spectrum: 637k eigenvalues exactly at $\gamma+\delta$ (the false bulk) + coupled satellites.}
\end{frame}

\begin{frame}{The 1-Bit Quantum Filter: the maths}
  \cn{[arXiv:2601.07766 --- Chiotopoulos, Nicotra, Scriven \emph{et al.}]}\\[2pt]
  Choose the evolution time $t=\pi/(\gamma+\delta)$. One-qubit QPE on $e^{iAt}$ leaves, per eigencomponent,
  \[ |u_k\rangle\Big(\cos\tfrac{\lambda_k t}{2}\,|0\rangle_{\rm clock} - i\,\sin\tfrac{\lambda_k t}{2}\,|1\rangle_{\rm clock}\Big), \]
  and the \texttt{X--CX--X} ancilla flip + post-selection keeps each mode with weight
  \[ \boxed{\;w_Q(\lambda)=\cos\Big(\frac{\pi}{2}\,\frac{\lambda}{\gamma+\delta}\Big)\;}\qquad
     w_Q(\gamma+\delta)=0\;\;\text{(the notch, placed \emph{exactly} on the false bulk)}. \]
  \centering\figh{0.4\textheight}{\SL/single_event_400/notch_rotation.png}\\
  \cn{The $0.25$ isolated-false bulk (637k segments) is rotated to $0$; the true band survives.}
\end{frame}

\begin{frame}{Why the 1BQF wins --- in brief}
  \begin{itemize}
    \item \textbf{One clock qubit} instead of an $n_{\rm time}$-qubit eigenvalue register.
    \item \textbf{Depth} cut by up to $10^4$ vs full HHL.
    \item \textbf{Sparse} $e^{iAt}$ (two-level/Givens decomposition) $\Rightarrow$ gate complexity $O(\sqrt{N}\log N)$.
    \item \textbf{Read-out solved}: the answer is binary (on/off the notch), no amplitude tomography.
    \item \textbf{NISQ-validated}: Quantinuum H2 and IBM Heron noise models.
    \item \textbf{Enables PV reconstruction} as classical post-processing.
  \end{itemize}
  \vspace{4pt}
  \centering\figh{0.3\textheight}{\PF/hhl_why_one_bit.png}
\end{frame}

\begin{frame}{The conceptual reformulation: inversion $\to$ filtering}
  \begin{columns}[T]
    \begin{column}{0.52\textwidth}
      \begin{itemize}
        \item We never needed the \emph{numbers} in $A^{-1}\mathbf{b}$ --- only a \textbf{binary decision} per segment: on or off.
        \item The spectrum is \textbf{bimodal}: false bulk \emph{at} $\gamma+\delta$, true band \emph{away} from it.
        \item So \textbf{one bit of spectral information classifies}: ``matrix inversion'' becomes ``\textbf{project out the known-bad eigenspace}''.
        \item HHL's eigenvalue \emph{precision} was wasted effort for this problem --- the 1BQF spends one bit exactly where the physics puts the information.
      \end{itemize}
    \end{column}
    \begin{column}{0.48\textwidth}
      \centering\fig{\textwidth}{\PF/eigenvalue_filter.png}
      \cn{$1/\lambda$ reweighting (HHL) vs the single-notch filter (1BQF).}
    \end{column}
  \end{columns}
\end{frame}

\begin{frame}{Every cluster type in segment space --- and its block of $A$}
  \centering
  \figw{0.97\textwidth}{\PF/A_blocks_segment_space.png}
\end{frame}

% ==================================================================
\section{Results: efficiency, false rate, and the trade-off}
\begin{frame}{Results: segment efficiency stays high; the false rate grows}
  \begin{columns}[T]
    \begin{column}{0.5\textwidth}
      \centering \cn{Classical} \\ \fig{\textwidth}{\SL/solver_segment_efficiency/classical_2x2_g3_base.png}
    \end{column}
    \begin{column}{0.5\textwidth}
      \centering \cn{1BQF at matched efficiency} \\ \fig{\textwidth}{\PF/quantum_2x2_matched_eff.png}
    \end{column}
  \end{columns}
  \vspace{-2pt}
  {\footnotesize \textbf{Both solvers: efficiency $\approx100\%$ across multiplicity; the false rate grows with $T$} (classical $0\!\to\!21\%$ by $T{=}1000$; quantum $0\!\to\!44\%$ by $T{=}400$) --- the coupled cross-track clusters.}
\end{frame}

\begin{frame}{The efficiency--false-rate trade-off}
  \begin{columns}[T]
    \begin{column}{0.5\textwidth}
      \centering\fig{\textwidth}{\SL/quantum_false_rate_investigation/operating_point_tradeoff.png}
    \end{column}
    \begin{column}{0.5\textwidth}
      \begin{itemize}
        \item Recovering \emph{every} true segment forces the threshold into the false band: \textbf{high efficiency and high false rate come together}.
        \item Quantum at matched efficiency ($\approx100\%$): false rate $\to 44\%$ at $T{=}400$.
        \item The alternative fixed-$\tau$ point (eff 75\%, far 1.7\%) only \emph{hides} the trade-off by discarding a quarter of the true segments.
        \item The \textbf{operating curve}, not a single point, is the result.
      \end{itemize}
    \end{column}
  \end{columns}
\end{frame}

\begin{frame}{Why the trade-off is fundamental}
  \begin{columns}[T]
    \begin{column}{0.52\textwidth}
      \begin{itemize}
        \item The notch removes only the \textbf{isolated} false class (it \emph{is} the bulk eigenvalue).
        \item Surviving false positives $=$ \textbf{coupled cross-track bridges/hubs}: \emph{good}, off-notch eigenvalues.
        \item In QPE phase they fall in the \textbf{same band as true segments} --- one zero cannot separate them.
        \item \textbf{Topological degeneracy}: a false chain shaped like a true track has an \emph{identical} $A$-block $\Rightarrow$ no threshold, classical or quantum, can split them.
      \end{itemize}
    \end{column}
    \begin{column}{0.48\textwidth}
      \centering\fig{\textwidth}{\BIF/phase_filter_map.png}
    \end{column}
  \end{columns}
\end{frame}

\begin{frame}{Class by class at matched efficiency ($T{=}400$)}
  \centering\figw{0.94\textwidth}{\PF/per_class_matched_T400.png}
  \vspace{-2pt}
  {\footnotesize The quantum keeps \textbf{all} true segments --- and with them the coupled false clusters: the price of efficiency.}
\end{frame}

% ==================================================================
\section{Attacking the survivors: the bifurcation term}
\begin{frame}{Add a bifurcation (fork) penalty to the Hamiltonian}
  The survivors are \textbf{competing continuations} --- segments fanning from a shared hit. Penalise them:
  \[ \boxed{\,A' = (\gamma+\delta)\,I - C + \beta B\,},\qquad \mathbf{b}=\delta\mathbf{1}, \]
  with $B$ the \textbf{fork adjacency}: segments sharing a start/end hit that are \emph{not} collinear continuations.
  Every segment touches exactly two hits $\Rightarrow$ the diagonal stays constant $\Rightarrow$ \textbf{still 1BQF-compatible}.
  \medskip
  \centering\figw{0.88\textwidth}{\PF/fork_matrix.png}
\end{frame}

\begin{frame}{The fork term: dense breaks the quantum --- $\varepsilon$-windowed fixes it}
  \begin{columns}[T]
    \begin{column}{0.5\textwidth}
      \centering\fig{\textwidth}{\BIF/eps_fork_quantum.png}
    \end{column}
    \begin{column}{0.5\textwidth}
      \centering\fig{\textwidth}{\BIF/eps_fork_sparsity.png}
    \end{column}
  \end{columns}
  \begin{itemize}
    \item \textbf{Dense} fork ($O(T^3)$ pairs): classically a benign uniform down-scaling (AUC $\approx1$) but \textbf{destroys the 1BQF} (AUC $\to0.55$) and the sparsity ($\sim$20$\times$ slower).
    \item \textbf{$\varepsilon$-windowed} $B_\varepsilon$ (only near-collinear competitors): \textbf{sparse}, \textbf{targeted} (false rate $\to0$ at $\sim$2\% efficiency cost), and \textbf{1BQF-safe} (AUC $0.96$--$0.99$).
  \end{itemize}
\end{frame}

% ==================================================================
\section{Next steps: QSVT}
\begin{frame}{Beyond one notch: Quantum Singular Value Transformation}
  \begin{columns}[T]
    \begin{column}{0.54\textwidth}
      \begin{itemize}
        \item The 1BQF is QSVT restricted to $p=\cos$ --- \textbf{one} notch.
        \item \textbf{QSVT} (Gily\'en--Su--Low--Wiebe) realises \emph{any} bounded polynomial $p(A)$ from a block-encoding of $A$ + $d$ phase rotations.
        \item Target: a \textbf{band-limited inverse} --- $p(\lambda)\approx1/\lambda$ on true-track eigenvalues, $p(\lambda)\approx0$ on every false mode (\emph{many} notches).
        \item Works \emph{iff} false and true eigenvalues are separated --- on real $T{=}200$ events the false modes sit $\ge0.2$ from every true eigenvalue.
      \end{itemize}
    \end{column}
    \begin{column}{0.46\textwidth}
      \centering\fig{\textwidth}{\QS/polynomial_design.png}
    \end{column}
  \end{columns}
\end{frame}

\begin{frame}{QSVT: a polynomial filter lifts true/false separation}
  \begin{columns}[T]
    \begin{column}{0.5\textwidth}
      \centering\fig{\textwidth}{\QS/efficacy_frontier.png}
    \end{column}
    \begin{column}{0.5\textwidth}
      \begin{itemize}
        \item A degree-$\sim$30 polynomial lifts true/false \textbf{AUC} from $0.49$ (classical) / $0.64$ (1-bit) to $\mathbf{0.88}$.
        \item Nulls \emph{all} off-length bridges and $\sim$95\% of hubs.
        \item \textbf{Cost:} $\sim$30 block-encoding calls $\times$ $\sim$2$\times$ amplitude-amplification $\approx$ tens$\times$ the 1-bit depth; sparsity preserved.
        \item \textbf{Irreducible floor:} a false bridge of the \emph{same length} as a true track (0\% at $T\le200$; grows with density).
      \end{itemize}
    \end{column}
  \end{columns}
\end{frame}

% ==================================================================
\section*{Summary}
\begin{frame}{Summary}
  \begin{itemize}
    \item Hits $\to$ segments $\to$ explicit kink-angle cost $C(\theta)$ $\to$ Ising-like $H(\mathbf{x})$ $\to$ stationarity $\Rightarrow$ \textbf{matrix inversion} $A\mathbf{x}=\delta\mathbf{1}$, $A=(\gamma+\delta)I-C$.
    \item \textbf{HHL} inverts it but needs deep circuits; \textbf{TrackHHL} restricts QPE to one bit ($10^4$ depth cut); the \textbf{1BQF} reframes inversion as a \textbf{single-notch spectral filter} --- $O(\sqrt N\log N)$, NISQ-ready.
    \item On the toy: \textbf{efficiency stays $\approx$100\%, the false rate grows with density} --- a fundamental trade-off set by topologically-degenerate cross-track clusters.
    \item Hamiltonian-level fix: the \textbf{$\varepsilon$-windowed fork term} (targeted, sparse, 1BQF-safe). Algorithm-level next step: \textbf{QSVT} polynomial filters (AUC $\to0.88$).
  \end{itemize}
  \vspace{4pt}
  \centering\cn{arXiv:2511.11458 (HHL) $\cdot$ arXiv:2601.07766 (1BQF) $\cdot$ toy: \texttt{LHCb\_VeLo\_Toy\_Model}}
\end{frame}

{\setbeamercolor{background canvas}{bg=thanksbg}
\begin{frame}[plain]\centering\vfill
  {\color{thanksfg}\Large\textbf{Thank you}}\\[4pt]
  {\color{thanksfg!85}\small Questions?}\\[16pt]
  \ThanksLogos
  \vfill
\end{frame}}

% ==================================================================
\appendix
\section*{Backup}
\begin{frame}{Backup: track-level metrics --- segments propagate to tracks}
  \centering\fig{0.62\textwidth}{\PF/track_metrics_vs_T.png}\\
  \cn{Wide $\varepsilon$ $\Rightarrow$ bridges merge tracks $\Rightarrow$ reconstruction efficiency collapses ($1.0\to0.42$); tight $\varepsilon$ stays perfect.}
\end{frame}

\begin{frame}{Backup: 1BQF--classical fidelity vs multiplicity}
  \centering\fig{0.62\textwidth}{\PF/cos_QC_vs_T.png}\\
  \cn{On the classical signal support the 1BQF is faithful: $\cos\theta_{QC}\approx0.97$, flat to $T{=}1000$.}
\end{frame}

\begin{frame}{Backup: fixed-threshold operating point ($\tau=0.35$)}
  \centering\fig{0.6\textwidth}{\SL/solver_segment_efficiency/quantum_2x2_g3_base.png}\\
  \cn{The same data at fixed $\tau$: eff 75\%, far 1.7\% --- the conservative end of the same operating curve.}
\end{frame}

\begin{frame}{Backup: the bad-eigenvalue ladder}
  \centering\fig{0.7\textwidth}{\SL/walkthrough_400/bad_eigenvalue_ladder.png}
\end{frame}

\begin{frame}{Backup: classical failure mode under noise}
  \centering\fig{0.85\textwidth}{\PF/noise_failure_mode.png}\\
  \cn{$\varepsilon$ scales with noise $\Rightarrow$ efficiency robust, but \textbf{purity} collapses under resolution smearing.}
\end{frame}

\begin{frame}{Backup: separability vs multiplicity}
  \centering\fig{0.7\textwidth}{\FRSE/separability_vs_T.png}\\
  \cn{Quantum starts cleaner, degrades faster; margin negative $\sim T{=}400$ --- the coupled-false ceiling.}
\end{frame}

\begin{frame}{Backup: $\gamma$ self-similarity of the notch}
  \centering\fig{0.85\textwidth}{\SL/quantum_false_rate_investigation/notch_universal_gamma.png}
\end{frame}

\begin{frame}{Backup: classical vs quantum per cluster (analytic)}
  \centering\fig{0.9\textwidth}{\SL/walkthrough_400/classical_vs_quantum.png}
\end{frame}

